\section{圧縮アルゴリズムの説明}
\subsection{全体像}
inputFileを開いて各文字の出現回数を記録し、その後出現回数が0の文字を取り除き、昇順でqueueのようなデータ構造で並ばせる。そのqueueの最も前、つまり出現回数が最も小さい二つを取り出し、Huffman Treeを作らせる。新たに生成されたrootをまたqueueに入れる。上記の動作をqueueに一つのデータしか残ってないまで繰り返すと、Huffman Treeが完全に生成される。生成された木の左と右を0と1に対応させ、各文字にcodeを付ける。更に、inputFileをrewind(inputFile)で先頭に戻し、もう一度読み直す。今度は各文字をその文字に付けられたcodeに置き換え、outputFileに書き込む。また、復元時のことを考えると、codeと各文字の対応関係、もしくは各文字の出現頻度を記録する必要がある。今回のプログラムでは後者を採用した。理由は後者の方がoutputFileが小さく、圧縮率を高められると思ったためだ。しかし、後者の場合、出現頻度からHuffman Treeを再構築する必要があり、やや解凍に時間かかるデメリットも存在する。
\subsection{個別の説明}
\subsubsection{FileHeaderの説明}
FileHeaderは各ファイルの名前、元のサイズと圧縮後のサイズ、圧縮された時間のような基本的なデータを記録している。それらは全てアーカイブに保存する。保存時はFile1のFileHeader、File1の内容、File2のFileHeader、File2の内容の順に保存する。理由はFileHeaderもしくはFileの内容をメモリで保存したらメモリの浪費に繋がるためファイル一つずつ書き込む。

\subsubsection{Struct  nodeの説明}
今回のプログラムでは、nodeはHuffman Treeの一つのnodeに対応する。char chはそのnodeの対応する文字であり、root、つまり他の二つのnodeにより作られた場合、対応する文字が無いため\verb|'\0'|に設定する。nextはnodeにqueneの構造を持たせるために設定したもので、最初はnodeとは別にstruct queneを作り、それにnode* dataとquene* nextを持たせようとしたが、quene\verb|->|data\verb|->|freqの参照と比べるのが面倒なのと、実際nodeに全部統一した方が効率もいいと思い、全部nodeに書き込んだ。また、nextとprev両方をnodeに持たせるのも可能だが、今回のプログラムではprevが無くても大して差は無さそうなためnextだけ持たせた。

\subsubsection{frequencyListの説明}
frequencyListのindexで各文字に対応させ、そのvalueはその文字が出現した回数に対応する。例えばA=65であるため、Aが出現するごとにfrequencyList[65]++を実行し、その出現を記録する

\subsubsection{insertNodeの説明}
newNodeと今の*headのfreqを比べ、newNodeが*headより前に並ぶ場合、headも更新する。それ以外の場合はnewNodeが今のポインター、つまりtempのnextのnodeの前か後ろのどちらに並ばせるべきかを判断する。ただし、tempのnextのnodeが存在しない場合は、tempの後ろにそのまま並ばせる
\subsubsection{compressFileToArchiveの説明}
このプログラムではrb,wbのような二進法でファイルを読み取るもしくは書き込む。この関数は特殊な場合、つまりファイルが空もしくは全部同じ文字である場合を考慮した。上記以外の場合は、fileSizeでファイルのサイズを保存し、それに応じてfileDataがメモリを確保し、一気にファイルを全部読み込む。frequencyListでfileDataの中身、つまり各文字の出現回数を記録する。codesがこのプログラムにおける唯一のグローバル変数であるため、先に全部初期化する。また、8bit = byteであるためcompressedSizeを計算する際、8で割った。もし小数が存在する場合、1を足す。上記が終わった後、FileHeaderにデータを書きこむ。最後に、FileHeaderのデータ、復元用のfrequencyListおよび肝心な圧縮後のデータをarchiveに書き込む。これでcompressの作業が終わり、必要なデータをprint outしてメモリをも解放し、return 1で終わる

